\section{Development Testing}

Testing during development consisted mostly of acceptance testing of individual features as they were implemented and regression testing to ensure that any new features did not break existing functionality \cite{anirbanSOFTWAREQUALITYASSURANCE2015}.

\subsection{Feature Testing}

Each component interaction was tested manually upon implementation to verify correct behavior.
This included:

\begin{itemize}
    \item \textbf{Grabbing and releasing:} Ensuring objects could be picked up and dropped reliably using the VR controllers;
    \item \textbf{Connection point detection:} Verifying that components would correctly identify valid mounting locations and snap into place when released within the acceptable tolerance;
    \item \textbf{Rotation mechanics:} Testing that both physical hand rotation and joystick-based fine adjustments worked as intended;
    \item \textbf{Audio and visual feedback:} Confirming that appropriate cues were triggered upon successful or failed placement attempts.
\end{itemize}

\subsection{Integration Testing}

After individual features were verified, integration testing was performed to ensure the complete PC building workflow could be executed from start to finish.
This involved building a complete PC configuration multiple times, verifying that:

\begin{itemize}
    \item Components could be installed in any order, but only valid connections could be made;
    \item The task list updated correctly as each step was completed;
    \item No component interactions interfered with previously placed parts.
\end{itemize}

\subsection{VR-Specific Considerations}

Given the nature of VR applications, additional testing focused on platform-specific concerns:

\begin{itemize}
    \item \textbf{Position translation:} Ensuring that real world hand movements accurately translated to in-app controller positions and to object manipulation;
    \item \textbf{Scale verification:} Ensuring component sizes matched real-world proportions to maintain immersion and educational accuracy during modelling.
\end{itemize}

\subsection{Limitations}

Due to time constraints and team size, formal unit testing and automated test suites were not implemented.
Testing was performed exclusively by the development team members, which may have introduced blind spots that were later identified during user testing (see Section \ref{sec:user-testing-results}).
Additionally, testing was mostly done on the team's PCs, with 